\documentclass[UTF8,a4paper,12pt]{ctexart}

\usepackage{amsmath}
\usepackage{amssymb}
\usepackage{geometry}
\usepackage{booktabs}
\usepackage{array}
\usepackage{listings}
\usepackage{xcolor}
\usepackage{hyperref}

\geometry{left=2.5cm,right=2.5cm,top=2.5cm,bottom=2.5cm}

\lstset{
    basicstyle=\ttfamily\small,
    keywordstyle=\color{blue},
    commentstyle=\color{gray},
    stringstyle=\color{teal},
    numbers=left,
    numberstyle=\tiny,
    breaklines=true,
    frame=single,
    columns=fullflexible
}

\title{Project 2：SR3 机械臂逆运动学解析推导与验证}
\author{}
\date{}

\begin{document}
\maketitle

\section{项目目标}

本项目围绕六自由度串联机械臂的逆运动学问题展开。逆运动学的目标是：
已知末端执行器相对于基坐标系的目标位姿
\[
{}^{0}T_{6}=
\begin{bmatrix}
{}^{0}R_{6} & {}^{0}p_{6}\\
0&1
\end{bmatrix},
\]
求出满足该位姿的关节变量
\[
q=[q_1,q_2,q_3,q_4,q_5,q_6].
\]

本报告以 SR3 机械臂为对象，基于 Project 1 中建立的 MDH 正运动学模型，
给出逆运动学求解思路。分析内容包括腕中心求解、前三关节位置子问题、
后三关节姿态子问题、多解性分析，以及 Matlab 代码测试结果。

\section{SR3 机械臂 MDH 参数}

本项目采用改进 DH（Modified DH, MDH）约定，单节变换矩阵定义为
\[
{}^{i-1}A_i
=R_x(\alpha_{i-1})T_x(a_{i-1})R_z(\theta_i)T_z(d_i).
\]

对应的齐次变换矩阵为
\[
{}^{i-1}A_i=
\begin{bmatrix}
\cos\theta_i & -\sin\theta_i & 0 & a_{i-1}\\
\sin\theta_i\cos\alpha_{i-1} &
\cos\theta_i\cos\alpha_{i-1} &
-\sin\alpha_{i-1} &
-d_i\sin\alpha_{i-1}\\
\sin\theta_i\sin\alpha_{i-1} &
\cos\theta_i\sin\alpha_{i-1} &
\cos\alpha_{i-1} &
d_i\cos\alpha_{i-1}\\
0&0&0&1
\end{bmatrix}.
\]

SR3 机械臂采用的 MDH 参数如下表所示，长度单位为 mm。

\begin{table}[h]
\centering
\begin{tabular}{ccccc}
\toprule
$i$ & $a_{i-1}$ & $\alpha_{i-1}$ & $d_i$ & $\theta_i$\\
\midrule
1 & 0   & $0$           & 344   & $q_1+\pi$\\
2 & 0   & $\frac{\pi}{2}$ & 0     & $q_2+\frac{\pi}{2}$\\
3 & 290 & $\pi$         & 0     & $q_3+\frac{\pi}{2}$\\
4 & 0   & $\frac{\pi}{2}$ & 290   & $q_4+\pi$\\
5 & 0   & $-\frac{\pi}{2}$& 0     & $q_5$\\
6 & 136 & $\frac{\pi}{2}$ & 103.5 & $q_6$\\
\bottomrule
\end{tabular}
\caption{SR3 机械臂 MDH 参数表}
\end{table}

因此在程序中关节变量与 MDH 角变量之间满足
\[
\theta_i=q_i+\mathrm{offset}_i,
\]
其中
\[
\mathrm{offset}
=
\left[
\pi,\frac{\pi}{2},\frac{\pi}{2},\pi,0,0
\right].
\]

求得 MDH 角变量后，需要再转换回真实关节变量：
\[
q_i=\theta_i-\mathrm{offset}_i.
\]

\section{逆运动学总体思路}

六自由度球腕机械臂的逆运动学通常可采用 Pieper 分离法。其核心思想是将
位姿求解拆成两个子问题：

\begin{enumerate}
    \item 位置子问题：由腕中心位置求解前 3 个关节变量；
    \item 姿态子问题：在前三关节确定后，由剩余姿态矩阵求解后 3 个关节变量。
\end{enumerate}

一般形式为
\[
{}^{0}T_6
={}^0A_1{}^1A_2{}^2A_3{}^3A_4{}^4A_5{}^5A_6.
\]

记
\[
{}^0R_6=R,\qquad {}^0p_6=p.
\]

若末端工具坐标系原点与球腕中心只有沿末端 $z_6$ 方向的偏置，则腕中心
可写为
\[
p_w=p-d_6R(:,3).
\]

但本项目所用 MDH 表中第 6 行同时存在
\[
a_5=136,\qquad d_6=103.5,
\]
即末端法兰存在横向偏置和轴向偏置。因此若直接采用
\[
p_w=p-d_6R(:,3)
\]
会遗漏 $a_5$ 对末端原点的影响。报告推导中应说明：严格求解时，应将末端
工具偏置完整剥离，或将第 6 轴偏置视为工具坐标系变换的一部分进行处理。

\section{腕中心的求解}

目标末端位姿记为
\[
T=
\begin{bmatrix}
R&p\\
0&1
\end{bmatrix}.
\]

若将第 6 坐标系中的固定工具偏置写成
\[
{}^6p_{\mathrm{tool}}
=
\begin{bmatrix}
a_5\\0\\d_6
\end{bmatrix},
\]
则末端原点可理解为腕部参考点经过该固定偏置后的位置。因此可近似写成
\[
p_w=p-R
\begin{bmatrix}
a_5\\0\\d_6
\end{bmatrix}.
\]

代入本机械臂参数，有
\[
p_w
=p-R
\begin{bmatrix}
136\\0\\103.5
\end{bmatrix}.
\]

需要注意的是，由于 MDH 坐标系中 $a_5$ 与 $d_6$ 的方向受腕部姿态影响，
在严格解析推导中，腕部偏置与 $\theta_5,\theta_6$ 存在耦合。
因此工程实现中采用正运动学反校验的方式剔除伪解，以保证最终解满足完整
MDH 模型。

\section{前三关节位置子问题}

前三关节主要决定腕中心位置。记腕中心为
\[
p_w=[x_w,y_w,z_w]^T.
\]

由腕中心在基坐标系 $xy$ 平面内的投影方向，可以得到第一关节角。
第一组肩部解为
\[
\theta_{1a}=\mathrm{atan2}(y_w,x_w).
\]

由于机械臂可从目标点另一侧到达，因此第二组肩部解为
\[
\theta_{1b}=\theta_{1a}+\pi.
\]

对任意一组 $\theta_1$，将腕中心投影到关节 2 和关节 3 构成的平面内。
定义
\[
r=x_w\cos\theta_1+y_w\sin\theta_1,
\]
\[
s=z_w-d_1.
\]

其中 $r$ 表示腕中心在旋转平面内的水平距离，$s$ 表示腕中心相对关节 2
平面的高度差。

对 SR3 机械臂，前臂和后臂的主要尺寸为
\[
L_1=290,\qquad L_2=290.
\]

在理想球腕模型下，位置子问题可化为平面二连杆问题：
\[
r=L_1\cos\theta_2+L_2\cos(\theta_2+\theta_3),
\]
\[
s=L_1\sin\theta_2+L_2\sin(\theta_2+\theta_3).
\]

由余弦定理可得
\[
\cos\theta_3
=
\frac{r^2+s^2-L_1^2-L_2^2}{2L_1L_2}.
\]

若
\[
|\cos\theta_3|>1,
\]
说明目标点超出机械臂工作空间，当前位置子问题无解。

当目标可达时，肘部有两组解：
\[
\theta_{3a}=\mathrm{atan2}
\left(
+\sqrt{1-\cos^2\theta_3},\cos\theta_3
\right),
\]
\[
\theta_{3b}=\mathrm{atan2}
\left(
-\sqrt{1-\cos^2\theta_3},\cos\theta_3
\right).
\]

对应地，
\[
\theta_2
=
\mathrm{atan2}(s,r)
-
\mathrm{atan2}
\left(
L_2\sin\theta_3,\,
L_1+L_2\cos\theta_3
\right).
\]

因此，前三关节位置子问题共有
\[
2\times 2=4
\]
组候选解，分别对应左肩/右肩和肘上/肘下。

\section{后三关节姿态子问题}

当前三关节确定后，可通过正运动学计算
\[
{}^0R_3.
\]

目标姿态满足
\[
{}^0R_6={}^0R_3{}^3R_6.
\]

因此可得腕部姿态矩阵
\[
{}^3R_6=({}^0R_3)^T{}^0R_6.
\]

记
\[
R_{36}={}^3R_6.
\]

对于后三关节，姿态矩阵可写为
\[
R_{36}
=R_x(\alpha_3)R_z(\theta_4)
 R_x(\alpha_4)R_z(\theta_5)
 R_x(\alpha_5)R_z(\theta_6).
\]

代入本 MDH 参数
\[
\alpha_3=\frac{\pi}{2},\qquad
\alpha_4=-\frac{\pi}{2},\qquad
\alpha_5=\frac{\pi}{2},
\]
得到
\[
R_{36}
=R_x\left(\frac{\pi}{2}\right)R_z(\theta_4)
R_x\left(-\frac{\pi}{2}\right)R_z(\theta_5)
R_x\left(\frac{\pi}{2}\right)R_z(\theta_6).
\]

根据矩阵元素可提取腕部角。令 $R_{36}=[r_{ij}]$，则有
\[
\cos\theta_5=-r_{23},
\]
\[
\sin\theta_5=\pm\sqrt{r_{13}^2+r_{33}^2}.
\]

因此
\[
\theta_5=\mathrm{atan2}(\sin\theta_5,\cos\theta_5).
\]

当
\[
|\sin\theta_5|>\varepsilon
\]
时，腕部非奇异。此时
\[
\theta_4=
\mathrm{atan2}
\left(
\frac{r_{33}}{\sin\theta_5},
\frac{r_{13}}{\sin\theta_5}
\right),
\]
\[
\theta_6=
\mathrm{atan2}
\left(
\frac{-r_{22}}{\sin\theta_5},
\frac{r_{21}}{\sin\theta_5}
\right).
\]

由于 $\sin\theta_5$ 可取正负两种符号，腕部姿态子问题存在两组翻转解。

当
\[
|\sin\theta_5|\approx 0
\]
时，机械臂处于腕部奇异位置，此时第 4 轴和第 6 轴的旋转效果发生耦合，
$\theta_4$ 和 $\theta_6$ 不再唯一。工程上可约定
\[
\theta_4=0,
\]
再由矩阵剩余元素求
\[
\theta_6=\mathrm{atan2}(-r_{12},r_{11}).
\]

\section{多解性分析}

一般 6R 球腕机械臂的逆运动学最多有 8 组解，其来源为：
\[
\text{肩部 2 解}\times
\text{肘部 2 解}\times
\text{腕部 2 解}
=8.
\]

具体而言：

\begin{itemize}
    \item 肩部多解：$\theta_1$ 可取相差 $\pi$ 的两组方向；
    \item 肘部多解：$\theta_3$ 的正负平方根对应肘上和肘下；
    \item 腕部多解：$\sin\theta_5$ 取正或负，对应腕部正翻和反翻。
\end{itemize}

在求得所有 MDH 角变量后，需要转换为真实关节角：
\[
\begin{aligned}
q_1&=\theta_1-\pi,\\
q_2&=\theta_2-\frac{\pi}{2},\\
q_3&=\theta_3-\frac{\pi}{2},\\
q_4&=\theta_4-\pi,\\
q_5&=\theta_5,\\
q_6&=\theta_6.
\end{aligned}
\]

同时还需要检查各关节是否满足机械臂关节限位：
\[
q_i^{\min}\le q_i\le q_i^{\max}.
\]

若存在多组满足约束的解，可根据参考关节角 $q_{\mathrm{ref}}$ 选择距离最近的一组：
\[
q^*
=
\arg\min_{q^{(k)}}
\left\|
\mathrm{wrapToPi}
\left(q^{(k)}-q_{\mathrm{ref}}\right)
\right\|_2.
\]

\section{代码实现说明}

Project 2 中主要文件如下：

\begin{itemize}
    \item \texttt{robot\_params.m}：保存 SR3 的 MDH 参数、offset 和关节限位；
    \item \texttt{forward\_kinematics.m}：根据 MDH 表计算正运动学；
    \item \texttt{inverse\_kinematics\_analytical.m}：实现逆运动学求解；
    \item \texttt{select\_ik\_solution.m}：按关节限位和参考角选择最优解；
    \item \texttt{jacobian\_geometric.m}：计算几何雅可比；
    \item \texttt{inverse\_kinematics\_numerical.m}：实现 DLS 数值逆运动学。
\end{itemize}

需要特别说明的是，本项目采用 MDH 约定。对 MDH 模型，关节 $i$ 的转轴并非
标准 DH 中直接取 $T_{i-1}$ 的 $z$ 轴，而是执行
\[
R_x(\alpha_{i-1})T_x(a_{i-1})
\]
之后、执行
\[
R_z(\theta_i)
\]
之前的 $z$ 轴。因此几何雅可比实现时需要按 MDH 结构重新确定每一列的关节轴。

\section{测试结果}

运行
\begin{lstlisting}[language=Matlab]
main_project2_test
\end{lstlisting}

得到测试结果如下：

\begin{lstlisting}
解析法 IK 自洽测试：5 / 5 PASS
解析法 IK 与 RTB 对比：3 / 3 PASS
几何雅可比测试：3 / 3 PASS
DLS 数值 IK 测试：3 / 3 PASS
\end{lstlisting}

其中，解析法 IK 自洽测试的基本流程为：

\begin{enumerate}
    \item 给定一组关节角 $q$；
    \item 通过正运动学计算目标位姿 $T$；
    \item 由逆运动学求得多组候选解 $Q_{\mathrm{all}}$；
    \item 将每一组候选解代入正运动学，验证是否重新得到目标位姿；
    \item 检查解集中是否存在与原始 $q$ 等价的解。
\end{enumerate}

测试结果表明，逆运动学求解得到的候选解均能通过正运动学反校验，
位置误差与姿态误差均满足阈值要求，说明所建立的 MDH 模型、逆运动学求解
过程和多解筛选策略是正确的。

\section{总结}

本文基于 SR3 机械臂的 MDH 模型，对六自由度机械臂逆运动学问题进行了分析。
首先由目标末端位姿分析腕中心位置，再将问题分解为前三关节的位置子问题和
后三关节的姿态子问题。前三关节存在肩部和肘部多解，后三关节存在腕部翻转
多解，因此一般情况下最多可得到 8 组逆运动学解。

由于当前 SR3 的 MDH 参数中包含末端法兰偏置，腕中心与末端原点之间并非只有
单一的 $d_6$ 轴向偏移。因此在代码实现中结合正运动学反校验和多初值搜索，
保证最终输出的候选解均严格满足完整 MDH 正运动学模型。测试结果显示，所有
设计用例均通过，证明逆运动学求解方法有效。

\end{document}
